\section{Full Algebraic Background}\label{app:min-algebra}

\subsection{State-Level Mergeability}

The classical mergeable-summary template has three moving parts: build a
state for each shard, merge states up the tree, and read the answer only from
the root state. It is the closed-form special case of the C-TreePO local-law
setup: instead of auditing residuals, the sketch proof shows directly that the
state still contains the information needed by the query. A
mergeable-summary scheme for query family $\mathcal{Q}$ is a triple
$(\mathrm{encode},\oplus,\mathrm{query})$ with
\[
\mathrm{encode}:\Strings\to\mathcal{S},\qquad
\oplus:\mathcal{S}\times\mathcal{S}\to\mathcal{S},\qquad
\mathrm{query}:\mathcal{S}\times\mathcal{Q}\to\mathcal{R}.
\]
Correctness asks that
\[
\mathrm{query}(\mathrm{encode}(u\concat v),q)
\approx
\mathrm{query}(\mathrm{encode}(u)\oplus\mathrm{encode}(v),q).
\]
The merge is on states, not on scalar child answers.
This is the common shape behind database algebraic aggregates, MUD
aggregation, and mergeable summaries
\citep{GrayEtAl1997DataCube,FeldmanEtAl2008MUD,Agarwal2013MergeableTODS}.
Associativity of $\oplus$ gives schedule invariance. Commutativity belongs to
unordered or symmetric tasks. Idempotence belongs to semilattice-like states,
not to mergeability in general.

\subsection{Ordered Text}

For ordered text the relevant algebra is a list, or free-monoid,
homomorphism:
\[
h(u\concat v)=h(u)\odot h(v)
\]
with associative but not necessarily commutative $\odot$
\citep{Gibbons1996ThirdHomomorphism}. The Markov boundary state is the
smallest example: the state contains first and last symbols, so swapping child
order changes the meaning. The Manifesto economic state is ordered for the
same reason: later qualifications, contrasts, and policy tradeoffs can change
how earlier commitments are read.

\subsection{State Versus Value}

The scalar task value need not merge. Distinct count is the canonical warning:
$|A|$ and $|B|$ do not determine $|A\cup B|$. HLL is mergeable because the
register state carries overlap-relevant information, and the readout is
applied after merge \citep{FlajoletEtAl2007}. In the play analogy, the sets
are character co-appearance pairs emitted by scenes; in the generic sketch
setting they are any stream items whose distinct union is the query. For
Manifesto economics, the seven-point score itself is not expected to merge
from child scores. The candidate mergeable object is the task-relative
evidence state.

\subsection{Qualified Equivalence}

The sketch-equivalence theorem is the state-level correspondence. Once the
state representation and readout are fixed, an oracle-respecting tree operator
$(g,f)$ satisfies C1/C2/C3 exactly when it factors through a mergeable scheme
$(\mathrm{encode},\oplus,\mathrm{query})$ on a state space $\mathcal{S}$. This
is not a claim that the scalar task values themselves merge; it is a claim
that the state carries enough information for the final readout.

\paragraph{Forward (mergeable $\Rightarrow$ local laws).}
Take a mergeable scheme with encode, merge, and query, and a decoder
$\mathrm{decode}:\mathcal{S}\to\Strings$. Define $g(x) =
\mathrm{decode}(\mathrm{encode}(x))$ on leaves and
$g(s_L\concat s_R) = \mathrm{decode}(\mathrm{encode}(s_L) \oplus \mathrm{encode}(s_R))$
on merges; let $f = \mathrm{query}\circ\mathrm{encode}$.
This is the whole nesting map. The summarizer $g$ is responsible for building
and merging state; the readout $f$ is applied only after that state has reached
the root. The assumptions needed are build validity, merge closure for the
union represented by concatenation, and readout correctness for every valid
state. C1 is the leaf case, C3 is the merge case, and C2 is needed only when
already-produced summaries can be summarized again. In Lean, this direction is
packaged as \path|local_laws_of_sketch|; the neural-operator subspace
version is \path|mergeableSketchSummaryClass_subset_exactLocalLawSubspace|.

\emph{C1 (sufficiency on leaves).} The hypothesis
\texttt{SketchLeafPreserving}
(\path|SketchSummaryOperators.lean|, line~190)
states that the encode--decode round-trip preserves $\fstar$ pointwise.
Applied to a leaf $b$, it gives
$\fstar(g(b))=\fstar(\mathrm{decode}(\mathrm{encode}(b)))=\fstar(b)$,
which is C1.

\emph{C2 (idempotence on stored summaries).} The hypothesis
\texttt{SketchSummaryCompatible} (line~203) gives
$\mathrm{decode}(\mathrm{encode}(s_L)\oplus\mathrm{encode}(s_R))=
g(\mathrm{decode}(s_L)\concat\mathrm{decode}(s_R))$, so re-encoding any
$s\in\operatorname{range}(g)$ lands in the same fiber. Combined with
the leaf hypothesis, $g(s)\sim_{\fstar}s$ holds on $\operatorname{range}(g)$,
which is C2.

\emph{C3 (merge consistency).} The hypothesis
\texttt{SketchMergeCompatible} (line~196) states that decoded child
sketches whose decoded values match $\fstar(x),\fstar(y)$ have a
merged decoded value matching $\fstar(x\concat y)$. Applied to
$s_L=\mathrm{encode}(u)$, $s_R=\mathrm{encode}(v)$, this gives
$\fstar(g(g(u)\concat g(v)))=\fstar(u\concat v)$, which is C3.

The packaging of these three checks into a \texttt{LocalLawsBundle}
is the formal theorem \path|local_laws_of_sketch|
(\texttt{SketchRecovery.lean}, line~40).

\paragraph{Reverse (local laws $\Rightarrow$ mergeable).}
Take $(g,f)$ satisfying C1, C2, C3. Set
$\mathcal{S}=\operatorname{range}(g)$ with $\mathrm{encode}=g$,
$\mathrm{decode}$ the identity inclusion $\mathcal{S}\hookrightarrow\Strings$,
$\oplus(s,t)=g(s\concat t)$, and $\mathrm{query}=f$. Then C1 gives
$\mathrm{query}\circ\mathrm{encode}=f\circ g\sim_{\fstar}\fstar$ on
leaves; C3 gives the mergeable correctness identity
$\mathrm{query}(\mathrm{encode}(u)\oplus\mathrm{encode}(v))=
f(g(g(u)\concat g(v)))\sim_{\fstar}\fstar(u\concat v)$; C2 is
automatic, since re-encoding an element of $\operatorname{range}(g)$
fixes its fiber by hypothesis.

For Agarwal-style summaries, the same state-level correspondence is relational
rather than canonical: validity is a relation \(V_\varepsilon(D,s)\), not
necessarily equality to one chosen state. The Lean names for that layer are
\begin{center}
\small
\path|RelationalMergeablePreferenceShape|\\
\path|RandomizedRelationalMergeablePreferenceShape|\\
\path|CTreePOToAgarwalTransform|
\end{center}

\paragraph{Approximate version.}
The exact equalities above relax to pointwise $\varepsilon$-budgets in
both directions. \texttt{SketchLeafApproxPreserving}
(\texttt{SketchRecovery.lean}, line~148) and
\texttt{SketchMergeApproxCompatible} (line~155) bound the violation
probabilities of the sketch-side hypotheses; these translate one-to-one
into the nodewise predicates \texttt{L1$\varepsilon$Node},
\texttt{L2$\varepsilon$Node} of \texttt{ApproximateLocalLaws.lean}, and
the bundle assembly is \texttt{approx\_bundle\_of\_sketch}
(\texttt{SketchRecovery.lean}, line~179). The audit reports the
empirical right-hand side of these inequalities on the realized calls
of the deployed tree, so the same equivalence carries the C-TreePO
recovery theorems
(\texttt{multi\_round\_proper\_of\_sketch},
\texttt{dpo\_equivalence\_of\_sketch},
\texttt{grpo\_equivalence\_via\_ZR\_of\_sketch})
to any mergeable sketch and to any audited learned summarizer.
